% Bagian: incompleteness
% Bab: representability-in-q
% Subbagian: c

\documentclass[../../../include/open-logic-section]{subfiles}

\begin{document}

\olfileid[id]{inc}{req}{cee}
\olsection{Fungsi-Fungsi dalam $C$}

Misalkan $C$ merupakan himpunan fungsi terkecil yang memuat
\begin{enumerate}
\item $0$,
\item penerus,
\item penjumlahan,
\item perkalian,
\item proyeksi, dan
\item fungsi karakteristik kesamaan, $\Char{=}$;
\end{enumerate}
serta tertutup terhadap
\begin{enumerate}
\item komposisi, dan
\item pencarian tak berbatas yang diterapkan pada fungsi reguler.
\end{enumerate}
Ingat bahwa pembatasan terakhir ini hanya berarti bahwa operasi~$\mu$ dapat
digunakan hanya ketika hasilnya total. Bandingkan hal ini dengan definisi
fungsi \emph{rekursif umum}: di sini kita telah menambahkan penjumlahan,
perkalian, dan $\Char{=}$, tetapi menghilangkan rekursi primitif.

Jelas bahwa segala sesuatu dalam $C$ bersifat rekursif, karena penjumlahan,
perkalian, dan $\Char{=}$ juga demikian. Kita akan menunjukkan bahwa
kebalikannya juga benar; ini sama dengan mengatakan bahwa, dengan unsur-unsur
lain dalam $C$, kita dapat menjalankan rekursi primitif.

\end{document}
